\chapter{Quatre problèmes de synthèse}

Chaque problème peut être travaillé en 25 à 40 minutes. Ce sont des sujets originaux d’entraînement, sans prétention à reproduire un format réglementaire d’examen.

\section*{Problème S1 -- Un jardin et son arrosage}

Un jardin rectangulaire mesure $18\unite m$ sur $12\unite m$. Un plan est réalisé au $1:200$. On y installe un bassin circulaire de rayon $2\unite m$.
\begin{enumerate}
\item Donner les dimensions du jardin et le rayon du bassin sur le plan.
\item Calculer l'aire de terrain restante, exactement puis au mètre carré près.
\item Une cuve cylindrique de rayon intérieur $0{,}7\unite m$ et de hauteur $1{,}2\unite m$ est remplie à $85\,\%$. Calculer son contenu en litres.
\item Un arrosage consomme $6\unite{L/m^2}$ sur toute la surface hors bassin. La cuve suffit-elle ? Conserver les valeurs exactes pour comparer.
\item Un tuyau doit suivre la diagonale du jardin. Calculer sa longueur exacte puis arrondie au décimètre.
\end{enumerate}

\section*{Problème S2 -- Un club et ses données}

Un club compare les offres $A(n)=30+4n$ et $B(n)=7n$ euros. Sur $20$ adhérents, les nombres de séances sont :
\[
\begin{array}{c|rrrr}\text{Séances}&5&8&10&15\\\hline\text{Effectif}&4&6&7&3\end{array}
\]
\begin{enumerate}
\item Déterminer quand les deux offres sont égales et quand A est strictement moins chère.
\item Calculer la moyenne, la médiane et l'étendue des nombres de séances.
\item Quel pourcentage d'adhérents a strictement intérêt à A ?
\item L'offre B augmente de $10\,\%$. Écrire la nouvelle fonction de coût et trouver le premier nombre entier de séances rendant A strictement moins chère.
\item Donner une formule de tableur, utilisant une référence de cellule pour $n$, qui calcule ce nouveau coût B.
\end{enumerate}

\section*{Problème S3 -- Une maquette et une hauteur}

Une maquette reproduit un bâtiment au $1:50$. Une face rectangulaire de la maquette mesure $16$ sur $10\unite{cm}$.
\begin{enumerate}
\item Donner les dimensions réelles de cette face et son aire.
\item Un volume de $240\unite{cm^3}$ sur la maquette correspond à quel volume réel en mètres cubes ?
\item Sur le terrain horizontal, un observateur situé à $12\unite m$ du bâtiment vise son sommet sous $32\degre$. Ses yeux sont à $1{,}65\unite m$ de hauteur. Calculer la hauteur visée.
\item Une autre mesure utilise un bâton de $1{,}5\unite m$ projetant une ombre de $2\unite m$, tandis que le bâtiment projette une ombre de $12\unite m$. Trouver la hauteur sous les hypothèses de similitude.
\item Les deux résultats coïncident-ils exactement ? Que faut-il contrôler avant d'accuser le théorème d'être faux ?
\end{enumerate}

\section*{Problème S4 -- Un jeu et son programme}

On tire successivement et avec remise deux boules dans une urne contenant $R_1,R_2,B_1,B_2,B_3$. Le joueur gagne si les deux couleurs sont identiques.
\begin{enumerate}
\item Compter les issues équiprobables et calculer la probabilité de gagner.
\item Calculer la probabilité de couleurs différentes par deux méthodes.
\item Décrire un algorithme simulant $5000$ parties et donnant la fréquence des gains.
\item Un programme choisit simplement rouge ou bleu avec une probabilité $1/2$ à chaque tirage. Pourquoi simule-t-il un autre jeu ?
\item Sans remise, quelle serait la probabilité de gagner ? Justifier par dénombrement.
\end{enumerate}